% SPDX-License-Identifier: LPPL-1.3c
% Copyright (C) 2026 Christophe Jorssen
% Author and Current Maintainer: Christophe Jorssen <christophe.jorssen@gmail.com>
% LPPL maintenance status: maintained
% This file is part of luafemm. See LICENSE and LICENSES.md for its terms.
\section{Problems, coordinates and physical paths}
\label{sec:model}
\begin{key}{/tikz/femm/problem=\marg{options} (default \normalfont empty)}
Create one model at the beginning of a |tikzpicture|. Put this key in the
picture options, not on a path or a nested scope. It installs |x=1mm,y=1mm|
and records the initial affine frame. There is one current model; nested
models inside the same picture are unsupported. Separate pictures create
independent models.
\end{key}

The following subkeys belong to |/femm/problem| and are normally supplied
inside |femm/problem={...}|. All lengths in this list are numbers in
\emph{model millimetres}, not graphic dimensions.

\begin{key}{/femm/problem/xmin=\meta{number} (initially -90)}
Left boundary of the rectangular air domain.
\end{key}
\begin{key}{/femm/problem/xmax=\meta{number} (initially 90)}
Right boundary; must exceed |xmin|. All material vertices must lie within
the domain, including its boundary.
\end{key}
\begin{key}{/femm/problem/ymin=\meta{number} (initially -75)}
Bottom boundary of the air domain.
\end{key}
\begin{key}{/femm/problem/ymax=\meta{number} (initially 95)}
Top boundary; must exceed |ymin|. By default, the potential is zero on the
entire outer rectangle. Use |\femmboundary| to change individual sides
(Section~\ref{sec:boundaries}). A finite exterior boundary is an approximation,
so enlarge the domain when testing the sensitivity of a result.
\end{key}
\begin{key}{/femm/problem/mesh size=\meta{positive number} (initially 2)}
Global target spacing of interior points and maximum length used to
subdivide input constraints. It is not a bound on every triangle diameter
or area. Smaller values increase both memory and solution time.
\end{key}
\begin{key}{/femm/problem/mesher=\meta{delaunay or grid} (initially delaunay)}
|delaunay| uses the constrained Lua triangulator and accepts general
polygonal interfaces. |grid| uses the structured algorithm,
which requires axis-aligned region edges. Local mesh sizes apply only to
|delaunay|. The standalone Lua constructor defaults to |grid|; the TikZ interface
explicitly selects |delaunay|.
\end{key}
\begin{key}{/femm/problem/mesh tolerance=\meta{nonnegative number} (initially 0)}
Vertex-merging distance. Zero selects $10^{-12}$ times the greater domain
dimension. A positive value must be below one percent of the global mesh
size. Features at this scale may merge. This tolerance controls geometric
construction, not the exact signs of the robust predicates.
\end{key}
\begin{key}{/femm/problem/curve tolerance=\meta{positive number} (initially .03)}
Default chord-distance tolerance for flattening material curves. It is
copied to the region options when the picture starts. It does not change
the original TikZ drawing and is distinct from mesh spacing. Measurement
profiles have a separate |curve tolerance| option.
\end{key}
\begin{key}{/femm/problem/minimum angle=\meta{number} (initially 0)}
Optional quality target in degrees, between 0 and 30 inclusive; zero disables
the angle criterion. Requires |mesher=delaunay| when nonzero. After constructing
the initial constrained mesh, split encroached segments and insert triangle
circumcentres until every angle meets the target, within floating-point audit
tolerance. Start with 20 degrees. A request that cannot be completed raises an
error; acute input corners may make the target impossible. See
Section~\ref{sec:mesh-quality} for the algorithm, limits and a complete example.
\end{key}
\begin{key}{/femm/problem/maximum area=\meta{nonnegative number} (initially 0)}
Optional upper bound on each triangle's area, in model millimetres squared.
Zero disables the area criterion. A nonzero value requires |mesher=delaunay|
and enables quality refinement even when |minimum angle=0|. Both criteria
apply together when enabled. This bound applies to the entire rectangular
domain, including background air; it is independent of local |mesh size|.
The Lua counterpart uses model-coordinate units squared, not SI units.
\end{key}
\begin{key}{/femm/problem/max Steiner points=\meta{nonnegative integer} (initially 5000)}
Maximum number of vertices added by quality refinement, counting both segment
midpoints and circumcentres. Initial constraint subdivision and lattice seeds
do not count towards this budget; the overall limit of 40000 mesh vertices
still applies. Zero permits no additional vertices. Exhausting the budget
raises an error, rather than accepting an under-refined mesh. The key has no
effect when both quality criteria are zero. Quality targets and this budget
are part of the persistent mesh descriptor; changes invalidate a saved mesh.
\end{key}
\begin{key}{/femm/problem/tolerance=\meta{positive number} (initially 1e-7)}
Target relative Newton residual, normalised by the initial residual. This
is an algebraic stopping criterion, not an estimate of the physical error.
\end{key}
\begin{key}{/femm/problem/max iterations=\meta{positive integer} (initially 60)}
Maximum number of Newton corrections. Nonconvergence raises an error;
the package does not silently use the last iterate.
\end{key}

\subsection{Declaring regions}
\begin{key}{/tikz/femm/region=\marg{options} (default \normalfont empty)}
Attach physical properties to one evaluated path with one or more closed
subcontours. Use |\path|,
|\draw| or |\filldraw| with |rectangle|, |circle|, |ellipse|, or |-- cycle|.
The capture happens after TikZ evaluates coordinates, transformations,
arcs, Bézier curves and rounded corners. Nodes on the path are annotations
and do not create additional regions. Stroke width does not add material.
\end{key}

\begin{key}{/femm/region/material=\meta{choice} (initially air)}
Select one of the 246 catalogue identifiers or a previously declared
custom material style. This is a PGF |.is choice| handler. Unknown values
produce an error, except that a declared |femm/materials/|\meta{name}
style is accepted as a custom choice. Catalogue identifiers are reserved.
\end{key}
\begin{key}{/femm/region/current density=\meta{expression or material} (initially material)}
Uniform signed $J_z$ in A/m$^2$. PGF expressions are accepted, including
|6000/(6*26)*1000000| for 6000 ampere-turns over a $6\times26$ mm$^2$
section. Positive current points out of the page. |material| inherits
the catalogue source; for most materials this is zero. Use an explicit
zero for a passive conductor. Currents are not renormalised if a path is
resized. The return section of a winding normally has the opposite sign.
\end{key}
\begin{key}{/femm/region/magnetization angle=\meta{expression} (initially 0)}
Prescribed direction in degrees from $+x$ towards $+y$, evaluated by PGF.
This is a direction in the model frame: rotating a local region path does
not automatically rotate its magnetisation. Rotating the entire picture
rotates both the rendered geometry and the displayed field.
\end{key}
\begin{key}{/femm/region/mesh size=\meta{nonnegative number} (initially 0)}
Local spacing inside the filled region and along every contour edge. Zero inherits the
global spacing; a positive value refines but does not coarsen that value.
Where refinement regions overlap, the smallest spacing wins. The region
still paints a material: declare a broad air refinement region before
the solids it contains, otherwise it replaces their material.
\end{key}
\begin{key}{/femm/region/curve tolerance=\meta{positive number}}
Override the picture's chord tolerance for this material path. The
initial value is copied from the problem at picture start.
\end{key}

\begin{keylist}[/tikz/femm]{material=\meta{choice},current density=\meta{expression},
  magnetization angle=\meta{expression},mesh size=\meta{number},curve tolerance=\meta{number}}
Aliases of the corresponding |/femm/region| subkeys. They can be inherited
from scopes or combined with graphics in a reusable TikZ style.
The alias |femm/mesh size| on a path means \emph{local} spacing; set the
global spacing inside |femm/problem|.
\end{keylist}

\begin{codeexample}[width=5cm]
\tikzset{passive/.style={
  femm/region={material=copper,
    current density=0},
  draw,fill=orange!25}}
\begin{tikzpicture}[femm/problem={
  xmin=-12,xmax=12,ymin=-12,ymax=12,
  mesh size=2}]
  \path[femm/region={material=air,
    mesh size=.6}]
    (-6,-5) rectangle (6,5);
  \filldraw[passive,rotate=15]
    (-4,-3) rectangle (4,3);
  \pic {femm mesh};
\end{tikzpicture}
\end{codeexample}

\subsection{Geometry and overlap rules}
A region may contain several individually simple closed contours.
Self-intersections, self-contact and backtracking within one contour are
errors. Native |even odd rule| and |nonzero rule| determine which portions
of a compound path receive its material; see Section~\ref{sec:domains}.
Distinct regions may overlap or intersect: \emph{the last declared
region whose fill contains a point determines material and current}.
The mesher splits all contour interfaces first, then classifies triangles
by their centroids. Air cavities remain part of the computed domain.

Simple node shapes with one closed background contour can also carry
|femm/region|. Their ordinary TikZ text sizing affects that contour.
For components whose dimensions must be independent of text, use the
node interfaces in Section~\ref{sec:components}. A pic may group several
physical paths, each with its own region declaration.

Use |femm/region| only on physical paths, not as |every path|. A clip
changes the visible part of a field or mesh, never the computed domain.
For profile paths, use |femm/profile| instead of |femm/region|.

The frame of the original picture is removed during capture and
reapplied to generated field and mesh graphics. Consequently |scale=2|
on the whole picture is a display change. A local scale inside the
picture changes the physical geometry. Fixed graphic dimensions, such
as the radius of |rounded corners|, retain TikZ's ordinary meaning.
Captured coordinates are rounded to $0.0001$ model unit, and points within
$0.001$ unit of a domain side are snapped to that side. Finer geometric
features cannot be recovered by reducing the mesh tolerance.

\subsection{Complete example: a nonlinear electromagnet}
This explicit-path example assembles a U core, its armature and the two
sections of a winding. Material and drawing options are collected into
TikZ styles. The displayed field and numeric diagnostics come from the
same physical paths; their declarations all precede the solve request.
\luafemmexample{u-electromagnet}

\subsection{Complete example: curved and transformed interfaces}
The next source rounds the core corners and tilts the armature with native
TikZ options. The mesh overlay exposes the flattened interfaces used by
the solver. Graphic clipping limits the displayed area without removing
any part of the computational domain.
\luafemmexample{curved-core}
